% =========================================================
\section{Loss-Corrected Quasi-One-Dimensional Model for Geometry Optimization}
\label{sec:loss_corrected_model}
% =========================================================

\subsection{Reduced-order performance closure}

\subsubsection{Need for a loss-corrected objective}

The ideal quasi-one-dimensional model developed in Section~\ref{sec:3} provides the core-flow
state, but it does not account for the loss of useful axial momentum caused by flow
divergence and viscous blockage. These effects are included through the momentum
efficiency

\begin{equation}
\boxed{
\eta_{mom}=\eta_{div}\eta_{BL}
},
\label{eq:momentum_efficiency}
\end{equation}

where \(\eta_{div}\) is the divergence efficiency and \(\eta_{BL}\) is the
boundary-layer velocity efficiency.

\subsubsection{Exit-flow divergence correction}

The divergence efficiency is evaluated from the exit-wall angle supplied by the
geometric model of Section~\ref{sec:geometry}. The axisymmetric approximation used by the simulator is

\begin{equation}
\boxed{
\eta_{div}
=
\frac{1+\cos\theta_{out}}{2}
}.
\label{eq:divergence_efficiency}
\end{equation}

This factor approaches unity as the exit flow becomes parallel to the nozzle axis and
decreases as \(\theta_{out}\) increases. The value of \(\theta_{out}\) is a derived
result of the contour formulation and is not an independent optimization variable.

\subsubsection{Boundary-layer correction}

The boundary-layer quantities are calculated with either the Quick screening model of
Section~\ref{sec:integral_boundary_layer} or the BLIMP-inspired model of
Section~\ref{sec:blimp_cebeci_smith}; their equations are not repeated here. The
blockage contribution enters through

\begin{equation}
\boxed{
\eta_{BL}
=
\frac{U_{e,eff}}{U_{e,ist}}
},
\qquad
0\leq\eta_{BL}\leq1,
\label{eq:boundary_layer_efficiency}
\end{equation}

where \(U_{e,eff}\) is the corrected exit velocity and \(U_{e,ist}\) is the
uncorrected quasi-one-dimensional value. Thus,
\(\eta_{BL}=1\) represents no predicted viscous velocity loss, whereas lower values
represent increasing boundary-layer blockage. The wall-shear contribution remains a
separate force coefficient, \(C_{F,fric}\), defined in
Eq.~\eqref{eq:blimp_friction_force}; it is not hidden inside \(\eta_{BL}\).

\subsubsection{Effective ambient thrust coefficient}

RocketCEA supplies the ideal momentum thrust coefficient
\(C_{F,mom}^{CEA}\) for each operating point and expansion ratio. The corrected
momentum contribution is

\begin{equation}
C_{F,mom}^{eff}
=
\eta_{div}\eta_{BL}C_{F,mom}^{CEA}.
\label{eq:corrected_momentum_cf}
\end{equation}

The ambient pressure contribution is evaluated separately:

\begin{equation}
C_{F,p}
=
\varepsilon\frac{p_e-p_a}{p_c}.
\label{eq:pressure_thrust_coefficient}
\end{equation}

The pressure term is positive for underexpanded operation, zero when \(p_e=p_a\), and
negative for overexpanded attached operation. It is not multiplied by the loss
efficiencies, since \(\eta_{div}\eta_{BL}\) is applied only to the momentum term. The
effective ambient thrust coefficient is therefore

\begin{equation}
\boxed{
C_{F,eff}
=
\eta_{div}\eta_{BL}C_{F,mom}^{CEA}
+
\varepsilon\frac{p_e-p_a}{p_c}
-C_{F,fric}
}.
\label{eq:effective_thrust_coefficient}
\end{equation}

The corresponding thrust estimate is

\begin{equation}
F_{eff}=C_{F,eff}p_cA_t.
\label{eq:effective_thrust}
\end{equation}

Since chamber pressure and throat area remain fixed during an optimization run,
maximizing \(C_{F,eff}\) is equivalent to maximizing \(F_{eff}\). The objective
therefore accounts simultaneously for divergence loss, boundary-layer blockage,
integrated wall friction and ambient pressure mismatch. The dependence of \(p_e\) on
\(\varepsilon\), including the ideal condition \(p_e=p_a\), follows from the
RocketCEA exit solution and Eq.~\eqref{eq:ideal_expansion}.

Ambient pressure is entered in the simulator in atmospheres and converted internally
according to

\begin{equation}
p_a[\mathrm{bar}]
=
1.01325\,p_a[\mathrm{atm}].
\label{eq:ambient_pressure_conversion}
\end{equation}

All pressures used in Eq.~\eqref{eq:effective_thrust_coefficient} are internally
converted to consistent units.

\subsubsection{Separated-flow handling and scope}

Equation~\eqref{eq:pressure_thrust_coefficient} assumes attached flow to the exit
plane. If RocketCEA classifies a candidate as separated, the candidate is assigned
zero fitness because the adopted exit-plane pressure term is no longer applicable.
The resulting loss correction remains a reduced-order engineering closure and must be
validated with higher-fidelity analysis or experimental data for final design use.
